\section{Questions at the Research Frontier}

\subsection{Between Statistical and Individual Fairness}
Given the limitations of extant notions of fairness, is there a way to get some of the ``best of both worlds''? In other words, constraints that are practically implementable without the need for making strong assumptions on the data or the knowledge of the algorithm designer, but which nevertheless provide more meaningful guarantees to individuals? Two recent papers, \cite{gerrymandering} and \cite{multical} (see also \cite{gerrymandering2,multical2} for empirical evaluations of the algorithms proposed in these papers), attempt to do this by asking for statistical fairness definitions to hold not just on a small number of protected groups, but on an exponential or infinite class of groups defined by some \emph{class of functions} of bounded complexity. This approach seems promising: because ultimately they are asking for statistical notions of fairness, the approaches proposed by these papers enjoy the benefits of statistical fairness: that no assumptions need be made about the data, nor is any external knowledge (like a fairness metric) needed.  It also better addresses concerns about ``intersectionality'', a term used to describe how different kinds of discrimination can compound and interact for individuals who fall at the intersection of several protected classes.

At the same time, the approach raises a number of additional questions: what function classes are reasonable, and once one is decided upon (e.g. conjunctions of protected attribures) what features should be ``protected''? Should these only be attributes that are sensitive on their own, like race and gender, or might attributes that are innocuous on their own correspond to groups we wish to protect once we consider their intersection with protected attributes (for example clothing styles intersected with race or gender)? Finally, this family of approaches significantly mitigates  some of the weaknesses of statistical notions of fairness by asking for the constraints to hold on average not just over a small number of coarsely defined groups, but over very finely defined groups as well. Ultimately, however, it inherits the weaknesses of statistical fairness as well, just on a more limited scale.

Another recent line of work aims to weaken the strongest assumption needed for the notion of individual fairness from \cite{awareness}: namely that the algorithm designer has perfect knowledge of a ``fairness metric''. \cite{KRR18} assume that the algorithm has access to an oracle which can return an unbiased estimator for the distance between two randomly drawn individuals according to an unknown fairness metric, and show how to use this to ensure a statistical notion of fairness related to \cite{gerrymandering,multical} which informally states that ``on average, individuals in two groups should be treated similarly if on average the individuals in the two groups are similar'' --- and this can be achieved with respect to an exponentially or infinitely large set of groups. Similarly, \cite{GJKR18} assumes the existence of an oracle which can identify fairness violations when they are made in an online setting, but cannot quantify the extent of the violation (with respect to the unknown metric). It is shown that when the metric is from a specific learnable family, this kind of feedback is sufficient to obtain an optimal regret bound to the best fair classifier while having only a bounded number of violations of the fairness metric. \cite{RY18} consider the case in which the metric is known, and show that a PAC-inspired approximate variant of metric fairness generalizes to new data drawn from the same underlying distribution. Ultimately, however, these approaches all assume that fairness is perfectly defined with respect to some metric, and that there is some sort of direct access to it. Can these approaches be generalized to a more ``agnostic'' setting, in which fairness feedback is given by human beings who may not be responding in a way that is consistent with any metric?


\subsection{Data Evolution and Dynamics of Fairness}
The vast majority of work in \emph{computer science} on algorithmic fairness has focused on one-shot classification tasks. But real algorithmic systems consist of many different components that are combined together, and operate in complex environments that are dynamically changing, sometimes because of the actions of the learning algorithm itself. For the field to progress, we need to understand the dynamics of fairness in more complex systems.

Perhaps the simplest aspect of dynamics that remains poorly understood is how and when components that may individually satisfy notions of fairness compose into larger constructs that still satisfy fairness guarantees. For example, if the bidders in an advertising auction individually are fair with respect to their bidding decisions, when will the allocation of advertisements be ``fair'', and when will it not? \cite{pipelines} and \cite{DI18} have made a preliminary foray in this direction. These papers embark on a systematic study of fairness under composition, and find that often the composition of multiple fair components will not satisfy any fairness constraint at all. Similarly, the individual components of a ``fair'' system may appear to be unfair in isolation. There are certain special settings, e.g. the ``filtering pipeline'' scenario of \cite{pipelines} --- modeling a scenario in which a job applicant is selected only if she is selected at every stage of the pipeline --- in which (multiplicative approximations of) statistical fairness notions compose in a well behaved way. But the high level message from these works is that our current notions of fairness compose poorly. Experience from differential privacy \cite{DMNS06,DR14} suggests that graceful degradation under \emph{composition} is key to designing complicated algorithms satisfying desirable statistical properties, because it allows algorithm design and analysis to be \emph{modular}. Thus, it seems important to find satisfying fairness definitions and richer frameworks that behave well under composition.

In dealing with socio-technical systems, it is also important to understand how algorithms dynamically effect their environment, and the incentives of human actors. For example, if the bar (for e.g. college admission) is lowered for a group of individuals, this might increase the average qualifications for this group over time because of at least two effects: a larger proportion of children in the next generation grow up in households with college educated parents (and the opportunities this provides), and the fact that a college education is achievable can incentivize effort to prepare academically. These kinds of effects are not considered when considering either statistical or individual notions of fairness in one-shot learning settings. The economics literature on affirmative action has long considered such effects --- although not with the specifics of machine learning in mind: see e.g. \cite{FV92,coates,Beck10}. More recently, there have been some preliminary attempts to model these kinds of effects in machine learning settings --- e.g. by modeling the environment as a markov decision process \cite{JJKMR17}, considering the equilibrium effects of imposing statistical definitions of fairness in a model of a labor market \cite{HC18}, specifying the functional relationship between classification outcomes and quality \cite{delayed}, or by considering the effect of a classifier on a downstream Bayesian decision maker \cite{KRZ18}. However, the specific predictions of most of the models of this sort are brittle to the specific modeling assumptions made --- they point to the need to consider long term dynamics, but do not provide robust guidance for how to navigate them. More work is needed here.

Finally, decision making is often distributed between a large number of actors who share different goals and do not necessarily coordinate. In settings like this, in which we do not have direct control over the decision making process, it is important to think about how to \emph{incentivize} rational agents to behave in a way that we view as fair. \cite{incentives} takes a preliminary stab at this task, showing how to incentivize a particular notion of individual fairness in a simple, stylized setting, using small monetary payments. But how should this work for other notions of fairness, and in more complex settings? Can this be done by controlling the flow of information, rather than by making monetary payments (monetary payments might be distasteful in various fairness-relevant settings)? More work is needed here as well. Finally, \cite{cd17} take a welfare maximization view of fairness in classification, and characterize the cost of imposing additional statistical fairness constraints as well. But this is done in a static environment. How would the conclusions change under a dynamic model?
\subsection{Modeling and Correcting Bias in the Data}

Fairness concerns typically surface precisely in settings where the available training data is already contaminated by bias.  The data itself is often a product of social and historical process that operated to the disadvantage of certain groups.   When trained in such data, off-the-shelf machine learning techniques may reproduce, reinforce, and potentially exacerbate existing biases.  Understanding how bias arises in the data, and how to correct for it, are fundamental challenges in the study of fairness in machine learning.

\cite{bolukbasi2016man} demonstrate how machine learning can reproduce biases in their analysis of the popular word2vec embedding trained on a corpus of Google News texts (parallel effects were independently discovered by \cite{CBN17}).  The authors show that the trained embedding exhibit female/male gender stereotypes, learning that ``doctor'' is more similar to man than to woman, along with analogies such as ``man is to computer programmer as woman is to homemaker''.  Even if such learned associations accurately reflect patterns in the source text corpus, their use in automated systems may exacerbate existing biases.   For instance, it might result in male applicants being ranked more highly than equally qualified female applicants in queries related to jobs that the embedding identifies as male-associated.

Similar risks arise whenever there is potential for feedback loops.  These are situations where the trained machine learning model informs decisions that then affect the data collected for future iterations of the training process.  \cite{lum2016predict} demonstrate how feedback loops might arise in predictive policing if arrest data were used to train the model\footnote{Predictive policing models are generally proprietary, and so it is not clear whether arrest data is used to train the model in any deployed system.}.  In a nutshell, since police are likely to make more arrests in more heavily policed areas, using arrest data to predict crime hotspots will disproportionately concentrate policing efforts on already over-policed communities.  Expanding on this analysis, \cite{ensign2018runawayfeedbackloops} find that incorporating community-driven data such as crime reporting helps to attenuate the biasing feedback effects. The authors also propose a strategy for accounting for feedback by adjusting arrest counts for policing intensity.  The success of the mitigation strategy of course depends on how well the simple theoretical model reflects the true relationships between crime intensity, policing, and arrests.  Problematically, such relationships are often  unknown, and are very difficult to infer from data.   This situation is by no means specific to predictive policing.

Correcting for data bias generally seems to require knowledge of how the measurement process is biased, or judgments about properties the data would satisfy in an ``unbiased'' world.  \cite{friedler2016possibility} formalize this as a disconnect between the \emph{observed space}---features that are observed in the data, such as SAT scores---and the unobservable \emph{construct space}---features that form the desired basis for decision making, such as intelligence.  Within this framework, data correction efforts attempt to undo the effects of biasing mechanisms that drive discrepancies between these spaces.  To the extent that the biasing mechanism cannot be inferred empirically, any correction effort must make explicit its underlying assumptions about this mechanism.  What precisely is being assumed about the construct space? When can the mapping between the construct space and the observed space be learned and inverted? What form of fairness does the correction promote, and at what cost?  The costs are often immediately realized, whereas the benefits are less tangible.  We will directly observe reductions in prediction accuracy, but any gains hinge on a belief that the observed world is not one we should seek to replicate accurately in the first place.  This is an area where tools from causality may offer a principled approach for drawing valid inference with respect to unobserved counterfactually `fair' worlds.

% \cite{bolukbasi2016man} propose debiasing algorithms that retain definitional associations to gender (e.g., queen is female, king is male) while removing gender associations to words that should be ``gender neutral''.


% Goodhart's law: Need to make sure what we're measuring reflects our goals.

% Causal models?
% In asking about the role that race or gender played in generating the observed data, we are asking a fundamentally causal question.  To the extent that this reduces the problem to one of causal discovery---determining the correct causal model based on observational data---this does not simplify the problem at all.  Causal discovery is often intractable.


%Observed differences reflect issues in the data or measurement process rather than meaningful patterns and trends that we would like to learn and replicate.

\subsection{Fair Representations}

Fair representation learning is a data de-biasing process that produces transformations (\textit{intermediate representations}) of the original data that retain as much of the task-relevant information as possible while removing information about sensitive or protected attributes.
This is one approach to transforming biased observational data in which group membership may be inferred from other features, to a construct space where protected attributes are statistically independent of other features.
First introduced in the work of \cite{zemel2013learning}, fair representation learning produces a de-biased data set that may in principle be used by other parties without any risk of disparate outcomes.  \cite{feldman2015certifying} and \cite{mcnamara2017provably} formalize this idea by showing how the disparate impact of a decision rule is bounded in terms of its balanced error rate as a predictor of the sensitive attribute.

Several recent papers have introduced new approaches for constructing fair representations.   \cite{feldman2015certifying} propose rank-preserving procedures for \textit{repairing} features to reduce or remove pairwise dependence with the protected attribute.  \cite{johndrow2017algorithm} build upon this work, introducing a likelihood-based approach that can additionally handle continuous protected attributes, discrete features, and which promotes joint independence between the transformed features and the protected attributes.  There is also a growing literature on using adversarial learning to achieve group fairness in the form of statistical parity or false positive/false negative rate balance \cite{edwards2015censoring,beutel2017data,zhang2018mitigating,madras2018learning}.

Existing theory shows that the fairness promoting benefits of fair representation learning rely critically on the extent to which existing associations between the transformed features and the protected characteristics are removed.  Adversarial downstream users may be able to recover protected attribute information if their models are more powerful than those used initially to obfuscate the data.  This presents a challenge both to the generators of fair representations as well as to auditors and regulators tasked with certifying that the resulting data is fair for use.
%Furthermore, as with other bias correction efforts, it isn't clear when removing all associations constitutes a reasonable fairness promoting strategy.  While this will ensure that the model predictions are independent of the protected attributes, such corrections may be going too far, or may be distorting the data in ways that lead to other forms of bias.
More work is needed to understand the implications of fair representation learning for promoting fairness in the real world.

% \begin{itemize}
%   \item Most existing approaches focus on removing associations.  \cite{johndrow2017algorithm}, \cite{madras2018learning}.
%   \item Recovering the protected feature permits disparate impact.  This creates a sort of arms race between regulators/auditors and modelers.
%   \item Causal approaches look at finer grained repairs, but rely on knowing or being able to estimate the causal structure.
%   \item Causal approach would model the full biasing mechanism. Probably intractable.
% \end{itemize}

% \begin{itemize}
%   \item Data generator needs to be more powerful than any downstream user
% \end{itemize}


%  What form of fairness does the correction promote, and at what cost?

\subsection{Beyond Classification}
Although the majority of the work on fairness in machine learning focuses on batch classification, batch classification is only one aspect of how machine learning is used. Much of machine learning --- e.g. online learning, bandit learning, and reinforcement learning --- focuses on \emph{dynamic} settings in which the actions of the algorithm feed back into the data it observes. These dynamic settings capture many problems for which fairness is a concern. For example, lending, criminal recidivism prediction, and sequential drug trials all are so-called \emph{bandit} learning problems, in which the algorithm cannot observe data corresponding to counterfactuals. We cannot see whether someone not granted a loan \emph{would have} paid it back. We cannot see whether an inmate not released on parole \emph{would have} gone on to commit another crime. We cannot see how a patient \emph{would have} responded to a different drug.

The theory of learning in bandit settings is well understood, and it is characterized by a need to trade off \emph{exploration} with \emph{exploitation}. Rather than always making a myopically optimal decision, when counter-factuals cannot be observed, it is necessary for algorithms to sometimes take actions that appear to be sub-optimal so as to gather more data. But in settings in which decisions correspond to individuals, this means sacrificing the well-being of a particular person for the potential benefit of future individuals. This can sometimes be unethical, and a source of unfairness \cite{explore}. Several recent papers explore this issue. For example, \cite{BBK17} and \cite{KMRWW18} give conditions under which linear learners need not explore at all in bandit settings, thereby allowing for best-effort service to each arriving individual, obviating the tension between ethical treatment of individuals and learning. \cite{RSVW18} show that the costs associated with exploration can be unfairly bourn by a structured sub-population, and that counter-intuitively, those costs can actually increase when they are included with a majority population, even though more data increases the rate of learning overall. However, these results are all preliminary: they are restricted to settings in which the learner is learning a linear policy, and \emph{the data really is governed by a linear model}. While illustrative, more work is needed to understand real-world learning in online settings, and the ethics of exploration.

There is also some work on fairness in machine learning in other settings --- for example, ranking \cite{YS17,CSV17}, selection \cite{selection,KR18}, personalization \cite{celis2017fair}, bandit learning \cite{infinitebandit,calibratedbandit}, human-classifier hybrid decision systems \cite{defer}, and reinforcement learning \cite{JJKMR17,DTB17}. But outside of classification, the literature is relatively sparse. This should be rectified, because there are interesting and important fairness issues that arise in other settings --- especially when there are combinatorial constraints on the set of individuals that can be selected for a task, or when there is a temporal aspect to learning.


%\ar{More topics?} 